\section*{الفصل \ref{ch:b1:observing-the-sky} --- التلسكوبات والكواشف والقياس الضوئي}

\begin{solution}{exo:b1:observing-the-sky:1}
فرق قدره 5 أقدار يعادل عاملًا قدره 100: فالنجم الذي قدره 1.0 يرسل 100 مرة من \omterm{def:b1:observing-the-sky:flux}{الدفق} أكثر. ريجل وبولاريس: $10^{0.4\times(2.02 - 0.13)} =
10^{0.756} = 5.70$؛ وريجل، ذو القدر الأصغر، أسطع.
\end{solution}

\begin{solution}{exo:b1:observing-the-sky:2}
$d = 1/0.19495 = \qty{5.13}{\parsec}$، أي $5.1295\times206265 =
\qty{1.058e6}{\au}$ و$5.1295\times3.2616 = \qty{16.7}{\lightyear}$.
\end{solution}

\begin{solution}{exo:b1:observing-the-sky:3}
$1.22\times\num{550e-9}/2.4 = \qty{2.80e-7}{rad} = \ang{;;0.058}$، أي نحو سبعة أضعاف دقّةً من رؤية فلكية قدرها \ang{;;0.40}: في الفضاء يبلغ التلسكوب حدّه الحيودي. أما \qty{6.5}{m} عند \qty{2.0}{\micro\metre}:
$1.22\times\num{2.0e-6}/6.5 = \qty{3.75e-7}{rad} = \ang{;;0.077}$.
\end{solution}

\begin{solution}{exo:b1:observing-the-sky:4}
$B - V$ يزداد من ريجل ($-0.03$) إلى ألتير (0.22)، وبولاريس (0.60)، وبيتلجوز (1.85): أحرّ النجوم ريجل، ثم ألتير فبولاريس فبيتلجوز. والبُعد يُخفت النطاقين بالعامل نفسه، فلا يغيّر $B - V$ (ما دام الغبار لا يحمّر الضوء في الطريق).
\end{solution}

\begin{solution}{exo:b1:observing-the-sky:5}
$X_1 = \sec\ang{30} = 1.1547$, $X_2 = \sec\ang{60} = 2.0000$.
$k = (6.590 - 6.412)/(2.0000 - 1.1547) = 0.178/0.8453 = \num{0.211}$ لكل كتلة هوائية، و$m_0 = 6.412 - 0.2106\times1.1547 = 6.169$.
\end{solution}

\begin{solution}{exo:b1:observing-the-sky:6}
$\mathrm{SNR} = \sqrt{400} = 20$. تنمو \omterm{def:b1:observing-the-sky:detector}{نسبة الإشارة إلى الضجيج} مع $\sqrt t$: من أجل 100،
$t = \qty{10}{s}\times(100/20)^2 = \qty{250}{s}$. ولخطأ قدره \num{0.01}
قدر تلزم $\mathrm{SNR} = 1.086/0.01 = 108.6$، إذن
$t = 10\times(108.6/20)^2 = \qty{295}{s}$.
\end{solution}

\begin{solution}{exo:b1:observing-the-sky:7}
البُعد هو $d = 1/0.00378 = \qty{264.6}{\parsec}$، \omterm{def:b1:observing-the-sky:absolute}{والقدر المطلق}
$M = 0.13 - 5\log_{10}26.46 = 0.13 - 7.11 = -6.98$. $L/L_\odot = 10^{0.4\times(4.83 + 6.98)} =
10^{4.724} = \num{5.3e4}$ في V.
\end{solution}

\begin{solution}{exo:b1:observing-the-sky:8}
$\pi\times1.2^2 = \qty{4.52}{m^2}$، و$25/4.52 = 5.5$: يجمع تلسكوب ويب العدد نفسه من الفوتونات في تعريض أقصر بمقدار 5.5 مرات. وقرص كامل قطره \qty{6.5}{m} مساحته $\pi\times3.25^2 = \qty{33.2}{m^2}$؛ والمرآة الحقيقية تمثّل منه $25/33.2 = \qty{75}{\percent}$، لأن القطع السداسية الثماني عشرة تترك حافة ذلك القرص غير مغطاة، وفجوات ما بين القطع والفتحة المركزية لا تجمع شيئًا.
\end{solution}

\begin{solution}{exo:b1:observing-the-sky:9}
(أ) $\qty{10}{\parsec} = \qty{3.0857e17}{m}$، و$F = \num{3.828e26}/(4\pi
\times(\num{3.0857e17})^2) = \num{3.828e26}/\num{1.1965e36} =
\qty{3.20e-10}{W.m^{-2}}$، و$1361/\num{3.20e-10} = \num{4.3e12}$ مرة أقل مما على الأرض.
(ب) $5\log_{10}(d/\qty{10}{\parsec}) = -26.74 - 4.83 = -31.57$، إذن
$d = \qty{10}{\parsec}\times10^{-6.314} = \qty{4.853e-6}{\parsec}$، أي
$\num{4.853e-6}\times206265 = \qty{1.001}{\au}$.
\end{solution}

\begin{solution}{exo:b1:observing-the-sky:10}
فوق الغلاف الجوي يكون \omterm{def:b1:observing-the-sky:flux}{الدفق} $F_0$؛ وتحته $F = F_0e^{-\tau_0X}$. إذن
$m - m_0 = -2.5\log_{10}(F/F_0) = 2.5\log_{10}e\;\tau_0X$، وهو $kX$ حيث
$k = 2.5\log_{10}e\;\tau_0 = 1.0857\,\tau_0$. ومن أجل $k = 0.160$:
$\tau_0 = 0.160/1.0857 = 0.147$. عند \omterm{def:b1:spherical-astronomy:celestial-sphere}{سمت الرأس} يصل $e^{-0.147} = 0.863$ من الضوء،
أي \qty{86}{\percent}؛ وعند $X = 2.31$، $e^{-0.147\times2.31} =
e^{-0.340} = 0.712$، أي \qty{71}{\percent}. وعند $z = \ang{85}$، $\sec z =
11.5$؛ لكن النموذج المتوازي المستويات يجعل المسار لا نهائيًا عند الأفق،
في حين يبقى المسار عبر غلاف منحنٍ منتهيًا، ويحني الانكسار الشعاع أيضًا:
قرب الأفق تبالغ الصيغة في تقدير \omterm{def:b1:observing-the-sky:extinction}{الكتلة الهوائية}، وتنمو أخطاؤها بسرعة.
\end{solution}

\begin{solution}{exo:b1:observing-the-sky:11}
(أ) $\mathrm{SNR} = Rt/\sqrt{(R + B)t}$، إذن $t = \mathrm{SNR}^2(R + B)/R^2
= 100\times2050/2500 = \qty{82}{s}$. (ب) تقسيم \omterm{def:b1:observing-the-sky:seeing}{الرؤية الفلكية} على اثنين يقسم
مساحة الصورة، ومن ثَمّ $B$، على 4: فيكون $B = 500$ و$t = 100\times550/2500 =
\qty{22}{s}$. (ج) عندما يكون $B \gg R$، يكون $t \approx \mathrm{SNR}^2B/R^2$. ويتزايد معدل المصدر $R$ مع $D^2$؛ أما معدل السماء $B$ فيتزايد مع المساحة المجمِّعة
$D^2$ مضروبةً في مساحة السماء في الصورة، $\theta^2$. إذن $t \propto
\theta^2D^2/D^4 = \theta^2/D^2$: فتقسيم \omterm{def:b1:observing-the-sky:seeing}{الرؤية الفلكية} على اثنين يعادل في نفعه مضاعفة القطر.
\end{solution}

\begin{solution}{exo:b1:observing-the-sky:12}
(أ) $N = 1200\times100 = \num{120000}$ خطًّا؛ وفي الرتبة الأولى تكون
قدرة التفريق \num{1.2e5}، و$\Delta\lambda = 550/(\num{1.2e5}) =
\qty{4.6e-3}{nm}$. (ب) $v = c\,\Delta\lambda/\lambda = c/(\num{1.2e5}) =
\qty{2.5}{km.s^{-1}}$. (ج) $\Delta\lambda = 550\times29.78/\num{299792} =
\qty{0.055}{nm}$، أي نحو 12 عنصر تفريق. وعلى الراصد أن يحسب،
لتاريخ الرصد ووقته، مركبة السرعة المدارية للأرض نحو النجم،
وأن يطرحها لتعطي سرعة النجم بالنسبة إلى الشمس.
\end{solution}

\begin{solution}{pb:b1:observing-the-sky:1}
\textbf{1.} $d = 1/0.37921 = \qty{2.637}{\parsec}$.

\textbf{2.} $2.637\times206265 = \qty{5.44e5}{\au}$;
$2.637\times\num{3.0857e16} = \qty{8.14e16}{m}$.

\textbf{3.} $2.637\times3.2616 = \qty{8.60}{\lightyear}$: غادر الضوء الذي نراه سيريوس قبل 8.6 سنة.

\textbf{4.} يتأرجح الاتجاه بمقدار $2p = \ang{;;0.758}$، أي نحو ضعف \omterm{def:b1:observing-the-sky:seeing}{الرؤية الفلكية}. ويُقاس التأرجح من موضع مركز الصورة، لا من حافتها: فبالمتوسط على فوتونات كثيرة يُحدَّد المركز بجزء صغير من عرض الصورة، وتزداد الدقة كلما عُدّت فوتونات أكثر.

\textbf{5.} $\Delta d/d = \Delta p/p = 0.001/0.37921 = \qty{0.26}{\percent}$،
أي $8.60\times0.0026 = \qty{0.023}{\lightyear}$.

\textbf{6.} $10^{0.4\times(0.03 + 1.46)} = 10^{0.596} = 3.94$.

\textbf{7.} $10^{0.4\times(26.74 - 1.46)} = 10^{10.112} = \num{1.29e10}$.

\textbf{8.} $10^{0.4\times(8.44 + 1.46)} = 10^{3.96} = \num{9120}$.

\textbf{9.} $z = \ang{90} - \ang{82.09} = \ang{7.91}$، $X = \sec\ang{7.91} =
1.0096$؛ والإخماد $0.15\times1.0096 = 0.151$ قدرًا؛ والقدر المقيس
$m = -1.46 + 0.151 = -1.31$.

\textbf{10.} $z = \ang{60}$، $X = 2.000$؛ والإخماد 0.30 قدرًا،
$m = -1.16$. والكسر النافذ هو $10^{-0.4\times0.30} = 0.759$:
\qty{24}{\percent} من الضوء يضيع.

\textbf{11.} عند الطلوع تكون $z = \ang{90}$ و$\sec z$ لا نهائية، بينما المسار الحقيقي عبر الغلاف المنحني طويل لكنه منتهٍ؛ ويحني الانكسار الشعاع أيضًا. قرب الأفق يكون الإخماد كبيرًا ويتغيّر بسرعة، والتصحيح المتوازي المستويات خاطئ.

\textbf{12.} $\num{3.0e4}\times10^{0.4\times(10.0 - 8.44)} = \num{3.0e4}
\times10^{0.624} = \num{3.0e4}\times4.21 = \num{1.26e5}$ إلكترونًا ضوئيًا في الثانية.

\textbf{13.} تلزم نسبة إشارة إلى ضجيج قدرها 1000 عدد $N = 1000^2 = \num{1.0e6}$
إلكترون ضوئي: $t = \num{1.0e6}/\num{1.26e5} = \qty{7.9}{s}$. والخطأ في القدر
هو $1.086/1000 = \num{0.0011}$.

\textbf{14.} $\num{3.0e4}\times10^{0.4\times(10.0 + 1.46)} = \num{3.0e4}
\times10^{4.584} = \num{1.15e9}$ إلكترونًا ضوئيًا في الثانية.

\textbf{15.} $0.20\times\num{1.15e9} = \num{2.30e8}$ في الثانية في أسطع بكسل،
وهو يمتلئ في $\num{1.0e5}/\num{2.30e8} =
\qty{4.3e-4}{s}$، أي أقل من نصف ملّي ثانية.

\textbf{16.} تعطي أجنحة سيريوس $\num{1.15e9}\times10^{-4} =
\num{1.15e5}$ في الثانية؛ ويعطي سيريوس~B $0.20\times\num{1.26e5} =
\num{2.52e4}$ في الثانية. فوهج سيريوس أسطع بمقدار 4.6 مرات من ذروة سيريوس~B.

\textbf{17.} يستقبل البكسل $\num{1.15e5} + \num{2.52e4} = \num{1.40e5}$
في الثانية، ويمتلئ في $\num{1.0e5}/\num{1.40e5} = \qty{0.71}{s}$. وفي
\qty{0.5}{s}: الإشارة $\num{2.52e4}\times0.5 = \num{1.26e4}$، والمجموع
$\num{7.01e4}$، والضجيج $\sqrt{\num{7.01e4}} = 265$، فتكون $\mathrm{SNR} =
\num{1.26e4}/265 = 48$. يمكن رصد سيريوس~B من حيث ضجيج الفوتونات؛ أما الصعوبة فهي أن الوهج يجب أن يُعرف ويُطرح بدقة تتجاوز بكثير خُمس قيمته، بينما تشبع نواة سيريوس في أقل من ملّي ثانية.

\textbf{18.} $m - M = 5\log_{10}(2.637/10) = 5\times(-0.579) = -2.89$.

\textbf{19.} $M = -1.46 + 2.89 = 1.43$.

\textbf{20.} $d = 1/0.13023 = \qty{7.679}{\parsec}$؛
$m - M = 5\log_{10}0.7679 = -0.57$، إذن $M = 0.03 + 0.57 = 0.60$.

\textbf{21.} فيغا، ذات \omterm{def:b1:observing-the-sky:absolute}{القدر المطلق} الأصغر، أكثر \omterm{def:b1:observing-the-sky:flux}{إشعاعية} بعامل قدره
$10^{0.4\times(1.43 - 0.60)} = 10^{0.332} = 2.15$. ولا يبدو سيريوس أسطع إلا لأنه أقرب.

\textbf{22.} $m = M_\odot + (m - M) = 4.83 - 2.89 = 1.94$: لكانت الشمس نجمًا ساطعًا إلى حدّ معقول يُرى بالعين المجردة.

\textbf{23.} $5\log_{10}(d/\qty{10}{\parsec}) = 6.0 - 1.43 = 4.57$، إذن
$d = 10\times10^{0.914} = \qty{82.0}{\parsec}$، أي نحو \qty{267}{\lightyear}.

\textbf{24.} $L/L_\odot = 10^{0.4\times(4.83 - 1.43)} = 10^{1.36} = 22.9$:
في النطاق V يفوق سيريوس ضوءَ \textbf{نحو 23 شمسًا}.
\end{solution}
